\documentclass[11pt]{article}
\usepackage[margin=1in]{geometry}
\usepackage{amsmath,amssymb,amsthm}
\usepackage{enumitem}

\newtheorem{theorem}{Theorem}
\newtheorem{proposition}[theorem]{Proposition}
\newtheorem{corollary}[theorem]{Corollary}
\newtheorem{remark}[theorem]{Remark}

\title{d5\_243: A periodic phase-carrier candidate family for the remaining graph-side search}
\author{}
\date{2026-03-22}

\begin{document}
\maketitle

\section{Purpose}
After the reverse reduction of \texttt{d5\_240}, the honest next question is not a
wide selector sweep but the following targeted realization problem:
\begin{quote}
Can the reduced omitted-row adjacent-transposition base, together with a
one-point cocycle defect, be realized by a small periodic local carrier?
\end{quote}
This note records the first clean answer.
There is an explicit periodic phase-carrier realization of the reduced target,
and on the checked small moduli no strictly smaller phase size survives inside
this family.
So the remaining graph-side search can be restricted further.

\section{The periodic phase-carrier family}
Fix an odd modulus $m>2$ and a positive integer $r$.
Let
\[
 p=(p_0,\dots,p_{r-1}),
 \qquad p_t\in \{\infty\}\cup \mathbb Z_m,
\]
be a row-pattern, and let $a\in\{+1,-1\}$.
Define
\[
 T_{m,r,a,p}\colon \mathbb Z_m^2\to \mathbb Z_m^2,
 \qquad (s,u)\mapsto (s+1,\tau_s(u)),
\]
where the row action is determined by the phase $t=s\bmod r$:
\begin{enumerate}[label=\textup{(\roman*)},leftmargin=2.7em]
\item if $p_t=\infty$, then $\tau_s$ is the identity;
\item if $p_t=\beta\in\mathbb Z_m$, then $\tau_s$ is the adjacent
      transposition swapping
      \[
      a s+\beta \quad\text{and}\quad a s+\beta+1
      \qquad (\bmod m).
      \]
\end{enumerate}
This is the natural periodic-carrier extension of the reduced family isolated in
\texttt{d5\_240}.
The carrier size is $r$.

\section{An explicit threshold candidate}
Write $m=2k+1$ and set
\[
 r=k+1=\frac{m+1}{2}.
\]
Consider the canonical phase pattern
\[
 p^{\mathrm{can}}=(\underbrace{0,\dots,0}_{k\text{ entries}},\infty).
\]
So phase classes $0,1,\dots,k-1$ use the unshifted moving adjacent transposition
and the last phase class is omitted.

\begin{proposition}[explicit periodic candidate]
\label{prop:explicit-periodic-candidate}
For $m=2k+1$ and $r=k+1$, the periodic phase-carrier map
\[
T_{m,r,+1,p^{\mathrm{can}}}
\]
is exactly the canonical omitted-row adjacent-transposition base from
\texttt{d5\_240} with defect row $k$.
In particular, it has a single orbit of size $m^2$ on $\mathbb Z_m^2$.
\end{proposition}

\begin{proof}
The rows are $s=0,1,\dots,2k$.
Their phase classes modulo $r=k+1$ are
\[
0,1,\dots,k,0,1,\dots,k-1.
\]
Hence rows $0,1,\dots,k-1$ use the adjacent transpositions
\[
(0\ 1),(1\ 2),\dots,(k-1\ k),
\]
row $k$ is omitted, and rows $k+1,\dots,2k$ use
\[
(k+1\ k+2),\dots,(2k\ 0).
\]
Therefore the row actions are exactly the moving adjacent transpositions
$(s\ \ s+1)$ for all rows $s\neq k$, with the single omitted row $s=k$.
That is precisely the omitted-row model from \texttt{d5\_240}.
The single-orbit conclusion follows from the base theorem proved there.
\end{proof}

\begin{corollary}[full grouped candidate]
\label{cor:grouped-candidate}
For the same explicit periodic candidate, a one-point cocycle defect yields a
single grouped orbit of size $m^3$.
\end{corollary}

\begin{proof}
By Proposition~\ref{prop:explicit-periodic-candidate}, the base has one orbit of
size $m^2$.
A one-point defect has total cocycle sum $1$, which is invertible modulo $m$.
So the single-cycle cocycle criterion from \texttt{d5\_240} applies.
\end{proof}

\section{Exact small-period pruning on the checked range}
The companion script
\texttt{d5\_242\_G1\_periodic\_phase\_carrier\_search.py}
performs an exact search over all phase patterns
\[
 p\in (\{\infty\}\cup\mathbb Z_m)^r,
 \qquad a\in\{\pm1\},
\]
for checked odd moduli
\[
 m=3,5,7,9,11,
\]
and for every strictly smaller phase size
\[
 1\le r\le \frac{m-1}{2}.
\]
It finds:
\begin{quote}
\emph{no single-orbit base candidates at all for any strictly smaller phase
size.}
\end{quote}
Equivalently, on the checked range, the periodic-carrier family has no honest
base realization below the threshold
\[
 r=\frac{m+1}{2}.
\]

\begin{remark}[what is exact here]
This is not a heuristic search.
For each checked pair $(m,r)$ the script exhausts the full finite family of all
patterns $p$ and both signs $a=\pm1$, and tests the full first-return
permutation on the row-zero coordinate.
Since the global map always advances $s$ by one, a single base orbit on
$\mathbb Z_m^2$ is equivalent to the first-return permutation being one
$m$-cycle.
\end{remark}

\section{What the first surviving layer looks like}
For the first surviving phase size
\[
 r=\frac{m+1}{2},
\]
the script also catalogs the exact survivor pool on the small moduli
$m=3,5,7,9$.
The observed counts are:
\[
6,\ 20,\ 84,\ 432.
\]
More importantly, every such checked survivor has:
\begin{enumerate}[label=\textup{(\arabic*)},leftmargin=2.8em]
\item exactly one omitted phase slot;
\item that omitted slot occurs at the final phase position in the canonical
      phase ordering used by the search.
\end{enumerate}
Representative examples are:
\[
[0,\infty],
\qquad
[0,0,\infty],
\qquad
[0,0,0,\infty],
\qquad
[0,0,0,0,\infty].
\]
So the explicit threshold candidate is not an isolated accident; it sits inside
an emerging exact survivor pool at the first possible carrier size.

\section{Consequence for the remaining graph-side search}
This still does \emph{not} close $G1$.
What it does do is reduce the live local-search question further:
\begin{enumerate}[label=\textup{(\arabic*)},leftmargin=2.8em]
\item any periodic realization attempt with carrier size strictly smaller than
      $(m+1)/2$ is already dead on the checked range;
\item there is a clean explicit threshold family that realizes the reduced
      omitted-row base and the one-point grouped defect exactly;
\item therefore the remaining live search should focus on how such a threshold
      carrier, or an equivalent non-periodic arithmetic carrier, is embedded
      into the actual selector-splice data behind $G1$.
\end{enumerate}
In particular, the graph-side search has moved beyond ``try another small local
rule grammar''. The honest remaining content is now a realization problem for a
carrier whose effective phase size grows with $m$, unless the manuscript splice
data encode an equivalent arithmetic coordinate more directly.

\section{Sanity checks recorded by the script}
The same companion script also checks the explicit threshold candidate on
\[
 m=3,5,7,9,11,13,15,17,19.
\]
For all these moduli it confirms:
\begin{enumerate}[label=\textup{(\arabic*)},leftmargin=2.8em]
\item the base orbit size is exactly $m^2$;
\item the first-return permutation on row $0$ is one $m$-cycle;
\item a one-point grouped cocycle defect gives one orbit of size $m^3$.
\end{enumerate}
So the explicit threshold candidate is not merely formal; it survives all
checked odd moduli in direct exact computation.

\end{document}
